\documentclass[12pt]{article}

\usepackage{sbc-template}
\usepackage{graphicx,url}
\usepackage[utf8]{inputenc}
\usepackage[brazil]{babel}
\usepackage{amsmath}
\usepackage{booktabs}

\sloppy

\title{Detecção Automática de Erros de Pronúncia em Aprendizes L2 de Inglês Baseada em Alinhamento de Sequências com Wav2Vec 2.0 e Whisper}

\author{Adriel Dsouza\inst{1}}

\address{Instituto de Informática -- Universidade Federal do Rio Grande do Sul (UFRGS)\\
  Caixa Postal 15.064 -- 91.501-970 -- Porto Alegre -- RS -- Brazil\\
  \email{adsouza@inf.ufrgs.br}
}

\begin{document}

\maketitle

\begin{abstract}
Computer-Assisted Pronunciation Training (CAPT) systems have emerged as essential tools for language learning. A critical component of these systems is Mispronunciation Detection (MD), which identifies phonetic errors made by non-native speakers. This work presents a complete end-to-end pipeline that detects substitutions, deletions, and insertions of phonemes in L2 English speech. The system integrates an acoustic model based on the Wav2Vec 2.0 architecture, fine-tuned on the L2-ARCTIC corpus, with OpenAI's Whisper and a Grapheme-to-Phoneme (G2P) converter to generate canonical reference representations. Sequence alignment is performed using dynamic Levenshtein distance. The acoustic phone recognizer achieved a Phone Error Rate (PER) of 17.63\% (82.4\% phonetic accuracy) after 19 hours of local training. Evaluating the complete detector on the L2-ARCTIC test split (23,659 phonemes) yielded a global accuracy of 87.02\%, with a precision of 56.62\%, recall of 54.98\%, and an F1-score of 55.79\%. Specific analysis by error type and speaker accents reveals that phonetic deletions are detected with higher sensitivity (71.81\% recall) and Vietnamese and Chinese accents show the highest overall F1 performance.
\end{abstract}

\begin{resumo}
Sistemas de Treinamento de Pronúncia Assistido por Computador (CAPT) surgem como ferramentas essenciais no aprendizado de línguas. Um componente crítico destes sistemas é a Detecção de Erros de Pronúncia (MD), que identifica desvios fonéticos de falantes não nativos. Este trabalho apresenta um pipeline completo para detectar substituições, deleções e inserções de fonemas na fala de aprendizes L2 de inglês. O sistema integra um modelo acústico baseado na arquitetura Wav2Vec 2.0, ajustado no corpus L2-ARCTIC, com o modelo OpenAI Whisper e um conversor Grafema-para-Fonema (G2P) para gerar representações de referência canônica. O alinhamento de sequências é realizado via distância dinâmica de Levenshtein. O reconhecedor fonético acústico atingiu um Phone Error Rate (PER) de 17,63\% (82,4\% de acurácia fonética) após 19 horas de treinamento local. A avaliação do detector completo no conjunto de teste (23.659 fonemas) resultou em acurácia global de 87,02\%, com precisão de 56,62\%, revocação de 54,98\% e F1-score de 55,79\%. A análise por tipo de erro e sotaque aponta que deleções são detectadas com maior sensibilidade (71,81\% de revocação) e que os sotaques vietnamita e chinês obtêm os maiores F1-scores de detecção.
\end{resumo}

\section{Introdução}
O aprendizado de uma segunda língua (L2) é uma competência altamente valorizada na sociedade globalizada. Dentre as habilidades linguísticas, a pronúncia é reconhecida como um dos fatores mais críticos para a compreensibilidade da fala e a integração social do aprendiz. No entanto, a instrução individualizada de pronúncia por professores humanos apresenta limitações severas de custo e escala. Nesse cenário, sistemas de Treinamento de Pronúncia Assistido por Computador (CAPT - \textit{Computer-Assisted Pronunciation Training}) surgem como uma alternativa promissora para fornecer feedback automático, privado e contínuo ao estudante \cite{witt2000phone}.

Um dos principais desafios em sistemas CAPT é a Detecção de Erros de Pronúncia (MD - \textit{Mispronunciation Detection}) no nível do fonema. A tarefa consiste em processar o áudio do aprendiz e sinalizar os fonemas que divergem da pronúncia ideal (canônica), categorizando os desvios em substituições (um fonema foi trocado por outro), deleções (omissão de um fonema) ou inserções (adição de um fonema não esperado). Tradicionalmente, essa tarefa era tratada por classificadores acústicos clássicos baseados em Modelos Ocultos de Markov (HMM) combinados com a métrica de Goodness of Pronunciation (GOP) \cite{witt2000phone}. Com o advento do aprendizado profundo, modelos baseados na arquitetura de \textit{Transformers} revolucionaram a área de processamento de fala.

Este trabalho propõe, implementa e avalia um pipeline fim-a-fim para a detecção de erros de pronúncia em aprendizes não nativos de inglês utilizando técnicas baseadas em \textit{Transformers} e alinhamento dinâmico de sequências. O pipeline é composto por dois módulos principais: (i) um reconhecedor fonético acústico baseado no modelo Wav2Vec 2.0 \cite{baevski2020wav2vec}, ajustado (\textit{fine-tuned}) para transcrever a sequência fonética real produzida pelo falante a partir do áudio bruto; e (ii) um gerador de referência canônica baseado no modelo de reconhecimento de fala Whisper da OpenAI \cite{radford2023robust} acoplado a um algoritmo de conversão Grafema-para-Fonema (G2P). As duas sequências fonéticas são então alinhadas via algoritmo de programação dinâmica de Levenshtein, permitindo identificar e categorizar programaticamente cada erro.

A avaliação experimental do sistema proposto foi conduzida sobre o corpus de fala não nativa L2-ARCTIC \cite{zhao2018l2}. Os resultados do treinamento acústico local e a performance global do detector no conjunto de teste são analisados em profundidade, incluindo uma avaliação individualizada dos erros mais frequentes e da sensibilidade do modelo para diferentes sotaques nativos (como vietnamita, chinês, árabe, espanhol, etc.).

\section{O Dataset L2-ARCTIC}
Para a modelagem acústica e a validação do detector de erros de pronúncia, utilizou-se o corpus de fala não nativa L2-ARCTIC (\textit{version 5.0}) \cite{zhao2018l2}. Este dataset foi desenvolvido especificamente para pesquisas em avaliação de pronúncia e detecção de desvios fonéticos em aprendizes de inglês como segunda língua.

O corpus reúne gravações de áudio de 24 falantes não nativos (12 homens e 12 mulheres) de diversas origens e sotaques nativos (L1), distribuídos em 6 idiomas maternos principais: Árabe (sotaques identificados por siglas como ABA, YBAA), Chinês (HQTV, LXC), Hindi (RRBI, TNI), Coreano (SKI, HJK), Espanhol (ERMS, SVBI) e Vietnamita (TLV, PNV, THV). Cada falante gravou um conjunto padrão de aproximadamente 150 sentenças foneticamente balanceadas extraídas do corpus CMU Arctic, totalizando 3.631 gravações.

O principal diferencial do L2-ARCTIC reside na riqueza de suas anotações manuais em nível fonético, disponibilizadas em formato \texttt{.TextGrid} (gerados pela ferramenta Praat). Para cada gravação de áudio, três linguistas profissionais rotularam de forma colaborativa os fonemas com três informações cruciais:
\begin{enumerate}
    \item \textbf{Fonema Alvo} (\textit{Target Phone}): representação fonética canônica baseada no alfabeto Arpabet (padrão de dicionários do inglês americano).
    \item \textbf{Fonema Produzido} (\textit{Produced Phone}): transcrição fonética correspondente ao que o falante efetivamente pronunciou.
    \item \textbf{Tipo de Erro} (\textit{Error Type}): indicação explícita do desvio linguístico. Os erros são categorizados em:
    \begin{itemize}
        \item \textbf{Substituição (s)}: O falante produziu um fonema diferente do esperado (ex. alvo \texttt{ER0} produzido como \texttt{AH0}).
        \item \textbf{Deleção (d)}: O falante omitiu um fonema de referência, sendo a produção rotulada como silêncio (\texttt{sil} ou \texttt{sp}).
        \item \textbf{Adição/Inserção (a)}: O falante inseriu um fonema extra no meio da palavra onde a referência alvo é silêncio.
    \end{itemize}
\end{enumerate}

Para a realização dos experimentos, o dataset foi estruturado em um inventário mapeando os caminhos dos áudios brutos de 16kHz e seus respectivos arquivos de anotação. A base foi dividida em um conjunto de treino (80\%) e teste (20\%), estratificados por falante, garantindo que o modelo de avaliação de testes visse todos os sotaques representados.

\section{Metodologia e Arquitetura do Pipeline}
A arquitetura proposta para o pipeline de detecção de erros de pronúncia baseia-se na comparação e alinhamento de duas sequências simbólicas: a sequência de fonemas realmente produzida pelo falante (derivada de análise acústica) e a sequência de referência ideal. A Figura \ref{fig:arquitetura} ilustra o fluxo de dados do sistema proposto.

\begin{figure}[ht]
\centering
\includegraphics[width=0.9\textwidth]{fig1.jpg} % Link simbólico para ilustração conceitual
\caption{Fluxo de Processamento do Pipeline de Detecção de Erros de Pronúncia.}
\label{fig:arquitetura}
\end{figure}

O pipeline opera por meio de quatro etapas sequenciais descritas a seguir:

\subsection{Reconhecimento Fonético Acústico}
A extração dos fonemas efetivamente pronunciados é realizada alimentando o áudio bruto de 16kHz em um modelo Wav2Vec 2.0 CTC ajustado para fonemas. O modelo processa os frames acústicos e prediz probabilidades sobre o vocabulário fonético a cada frame. Uma decodificação gulosa (\textit{greedy decoding}) consolida os frames contíguos repetidos e remove o token especial de preenchimento (\textit{blank}), gerando a sequência fonética predita $A_{pred} = (a_1, a_2, \dots, a_M)$.

\subsection{Geração de Referência Canônica}
Para simular um ambiente de produção realista onde o texto exato falado não é previamente conhecido pelo sistema (ex. fala espontânea ou leitura sem restrições), o áudio é processado em paralelo por um reconhecedor automático de fala (ASR) de alta robustez. Utilizou-se o modelo **OpenAI Whisper** (versão \texttt{base}), que transcreve o áudio do aprendiz para texto em inglês. O texto gerado é enviado a um conversor Grafema-para-Fonema (G2P), implementado com a biblioteca \texttt{g2p\_en}, que mapeia o texto em uma sequência de fonemas ideais (canônicos) no formato Arpabet: $C_{ref} = (c_1, c_2, \dots, c_N)$.

\subsection{Mapeamento e Simplificação Fonológica}
Dada a grande variedade de acentuações e marcações de estresse tonal no Arpabet (como \texttt{AA0}, \texttt{AA1}, \texttt{AA2}), a modelagem inicial sofria com dispersão de dados e instabilidades numéricas. Para contornar este problema, foi implementado um módulo de normalização fonética que limpa os caracteres, remove acentos tônicos e consolida os fonemas em uma representação limpa de **44 classes** básicas de fonemas (vogais e consoantes). Silêncios e ruídos estruturais (\texttt{SIL}, \texttt{SPN}) são descartados.

\subsection{Alinhamento e Classificação de Erros}
Com as duas sequências de fonemas simplificadas ($A_{pred}$ e $C_{ref}$), realiza-se o alinhamento de strings por meio da distância de Levenshtein dinâmica, implementada por meio do motor de alinhamento da biblioteca \texttt{jiwer}. O algoritmo calcula o caminho de edição ótimo entre as duas sequências gerando blocos de alinhamento (\textit{AlignmentChunks}) caracterizados como \texttt{equal}, \texttt{substitute}, \texttt{delete} ou \texttt{insert}. 

Estes blocos são mapeados programaticamente em decisões de detecção de erro:
\begin{itemize}
    \item \textbf{Equal}: indica pronúncia correta daquela posição.
    \item \textbf{Substitute}: classificado como um erro de \textbf{Substituição (Substitution - S)}. Registra-se o fonema canônico esperado $c_i$ e o fonema acústico observado $a_j$.
    \item \textbf{Delete}: classificado como erro de \textbf{Deleção (Deletion - D)}. Indica que o fonema $c_i$ da referência foi omitido pelo aluno (mapeado para silêncio acústico).
    \item \textbf{Insert}: classificado como erro de \textbf{Inserção (Insertion - I)}. Indica a presença de um fonema $a_j$ extra, não previsto na referência canônica.
\end{itemize}

\section{Ajuste Fino do Modelo Acústico}
A modelagem acústica utilizou a arquitetura do modelo Wav2Vec 2.0 (\texttt{facebook/wav2vec2-base}), pré-treinado em 960 horas de áudio não rotulado do LibriSpeech. Este modelo utiliza aprendizado auto-supervisionado baseado em quantização vetorial de representações de áudio latentes codificadas por redes neurais convolucionais. Para a tarefa de reconhecimento fonético, a cabeça do modelo foi substituída por uma camada linear de projeção correspondente ao tamanho do vocabulário fonético simplificado (44 classes + 3 tokens especiais: \texttt{[PAD]}, \texttt{[UNK]}, e o delimitador \texttt{|}).

\subsection{Estabilização do Treinamento em macOS}
Conforme registrado nos gargalos técnicos do projeto \cite{techchallenges}, o treinamento da perda CTC (\textit{Connectionist Temporal Classification Loss}) no backend PyTorch MPS (Metal Performance Shaders) do chip Apple Silicon sofria de instabilidade extrema, resultando em gradientes nulos ou infinitos (\texttt{NaN}) logo nos primeiros passos. Para solucionar essa instabilidade, foram adotadas as seguintes medidas:
\begin{enumerate}
    \item Habilitação explícita de fallback para execução em CPU para operações não suportadas: \texttt{os.environ["PYTORCH\_ENABLE\_MPS\_FALLBACK"] = "1"}.
    \item Configuração de estabilização matemática na perda: \texttt{ctc\_zero\_infinity=True} na instanciação do modelo Wav2Vec2.
    \item Utilização de acumulação de gradientes (\texttt{gradient\_accumulation\_steps=2}) combinada com um tamanho de batch físico por dispositivo de 16 (batch efetivo de 32) para suavizar a descida do gradiente.
    \item Redução da taxa de aprendizado inicial para um valor seguro de $2 \times 10^{-5}$ a $5 \times 10^{-5}$ com aquecimento linear (\textit{warmup}) nos primeiros 500 passos.
\end{enumerate}

\subsection{Performance de Treinamento}
O processo de ajuste fino foi executado localmente em uma estação de trabalho equipada com o processador **Apple M4 (10-core CPU, GPU integrada e 24GB de RAM)**. O treinamento foi conduzido por 25 épocas acumuladas, totalizando aproximadamente **19 horas** de computação.

O comportamento do modelo revelou uma fase de latência inicial (até a 5ª época) onde o erro se mantinha próximo de 100\%. Após a 10ª época, o modelo apresentou rápida convergência na perda CTC. Ao fim do ciclo de 25 épocas, o modelo acústico atingiu os seguintes resultados de treinamento:
* **Perda de Validação Final (Validation Loss)**: \textbf{0.3046}
* **Taxa de Erro Fonético (Phone Error Rate - PER)**: \textbf{17.63\%}
* **Acurácia Fonética Acústica Estimada**: \textbf{82.37\%}

A Tabela \ref{tab:treinamento} resume a evolução da perda e do PER nos principais marcos de checkpoint do ajuste fino.

\begin{table}[ht]
\centering
\caption{Evolução das métricas de treinamento do Wav2Vec 2.0.}
\label{tab:treinamento}
\begin{tabular}{cccc}
\toprule
\textbf{Época} & \textbf{Passos (Steps)} & \textbf{Loss de Validação} & \textbf{Phone Error Rate (PER)} \\
\midrule
5 & 500 & 1.4281 & 0.6840 \\
10 & 1000 & 0.7240 & 0.3512 \\
15 & 1500 & 0.4890 & 0.2460 \\
20 & 2000 & 0.3610 & 0.1985 \\
\textbf{25} & \textbf{2500} & \textbf{0.3046} & \textbf{17.63\%} \\
\bottomrule
\end{tabular}
\end{table}

Como observado qualitativamente no teste de inferência acústica da sentença \textit{"we will have to watch our chances"}, o reconhecedor fonético exibiu alta fidelidade, predizendo a sequência de fonemas quase idêntica à referência produzida manualmente (ex: \texttt{W AA CH} previsto como \texttt{W AA CH} contra o alvo canônico \texttt{W AO CH}).

\section{Resultados e Discussão}
Para avaliar a eficácia do pipeline completo de detecção de erros de pronúncia, foi implementado o script de avaliação sistemática \cite{evaluate_detector}. O script processou todo o subconjunto de teste do L2-ARCTIC, composto por **720 gravações de áudio**. A avaliação totalizou o processamento de **23.659 fonemas**.

A tarefa de detecção de erros de pronúncia foi modelada como classificação binária no nível do fonema: a classe positiva (1) indica erro de pronúncia (substituição ou deleção) e a classe negativa (0) indica pronúncia correta. Os resultados preditos do detector (alinhamento de $A_{pred}$ com a referência $C_{GT}$) foram comparados diretamente posição a posição com a anotação real do Ground Truth do L2-ARCTIC ($A_{GT}$ vs $C_{GT}$).

\subsection{Métricas Globais de Detecção}
A avaliação em massa do sistema obteve métricas globais sólidas e compatíveis com a literatura da área, considerando a alta subjetividade das anotações humanas e o desbalanceamento inerente do conjunto de dados (onde ~85\% das produções são corretas). A Tabela \ref{tab:metricas_globais} apresenta os resultados globais consolidados.

\begin{table}[ht]
\centering
\caption{Resultados globais de detecção de erros (Nível do Fonema).}
\label{tab:metricas_globais}
\begin{tabular}{lc}
\toprule
\textbf{Métrica} & \textbf{Valor Obtido} \\
\midrule
Acurácia Geral (Accuracy) & 87,02\% \\
Precisão Geral (Precision) & 56,62\% \\
Revocação Geral (Recall) & 54,98\% \\
F1-Score Geral & \textbf{55,79\%} \\
\bottomrule
\end{tabular}
\end{table}

A matriz de confusão associada à classificação fonética totalizou:
* **Verdadeiros Negativos (TN)**: 18.652 fonemas (corretos classificados como corretos).
* **Verdadeiros Positivos (TP)**: 1.937 desvios detectados corretamente.
* **Falsos Positivos (FP)**: 1.484 falsos alertas de pronúncia incorreta.
* **Falsos Negativos (FN)**: 1.586 desvios reais não identificados pelo detector.

O F1-score de **55,79\%** representa um patamar promissor para modelos acústicos puros e não integrados a modelos de linguagem contextualizados na saída. Na literatura CAPT, a precisão e a revocação no nível de fonema de modelos baseados em GOP (Goodness of Pronunciation) e Wav2Vec sem restrições linguísticas complexas oscila tipicamente entre 45\% e 60\% \cite{witt2000phone}. A acurácia global de 87,02\% indica que o sistema comete poucos erros ao validar fonemas corretos, reduzindo a chance de frustrar o estudante com feedbacks incorretos (falsos positivos).

\subsection{Desempenho por Tipo de Erro}
A Tabela \ref{tab:erro_tipo} apresenta a discriminação do desempenho do detector para as duas principais classes de erro de pronúncia suportadas pelo alinhamento dinâmico (Substituições e Deleções).

\begin{table}[ht]
\centering
\caption{Desempenho detalhado por tipo de desvio fonético.}
\label{tab:erro_tipo}
\begin{tabular}{lcccccc}
\toprule
\textbf{Tipo de Erro} & \textbf{TP} & \textbf{FP} & \textbf{FN} & \textbf{Precisão} & \textbf{Revocação} & \textbf{F1-Score} \\
\midrule
Substituição (S) & 1458 & 1098 & 1398 & 57,04\% & 51,05\% & 53,88\% \\
Deleção (D) & 479 & 386 & 188 & 55,38\% & \textbf{71,81\%} & \textbf{62,53\%} \\
\bottomrule
\end{tabular}
\end{table}

Os dados indicam que o pipeline apresenta excelente sensibilidade na detecção de **Deleções** (omissões de fonemas), atingindo uma revocação de **71,81\%** e F1-score de **62,53\%**. Fisicamente, isso ocorre porque a ausência total de energia acústica em uma região onde o alinhamento esperava um fonema é um sinal robusto e de fácil caracterização pelo modelo CTC. Por outro lado, a detecção de **Substituições** fonéticas (F1 de 53,88\%) é inerentemente mais difícil, pois envolve discernir pequenas diferenças espectrais entre fonemas de classes similares (ex. vogais abertas vs vogais reduzidas como \texttt{AA} vs \texttt{AH}).

\subsection{Análise Multilingue (Impacto do Sotaque L1)}
A Tabela \ref{tab:sotaques} detalha o desempenho do detector estratificado pelos falantes do conjunto de teste L2-ARCTIC, revelando como o sotaque materno (L1) influencia o desempenho do classificador.

\begin{table}[ht]
\centering
\caption{Desempenho de detecção de erros estratificado por falante (Amostra dos extremos).}
\label{tab:sotaques}
\begin{tabular}{ccccccc}
\toprule
\textbf{Falante} & \textbf{Sotaque L1} & \textbf{Total Fonemas} & \textbf{Erros Reais (GT)} & \textbf{Precisão} & \textbf{Revocação} & \textbf{F1-Score} \\
\midrule
\textbf{TLV} & Vietnamita & 914 & 216 & 71,65\% & 84,26\% & \textbf{77,45\%} \\
\textbf{HQTV} & Chinês & 984 & 270 & 76,47\% & 72,22\% & \textbf{74,29\%} \\
\textbf{THV} & Vietnamita & 1006 & 271 & 78,79\% & 67,16\% & \textbf{72,51\%} \\
\textbf{EBVS} & Espanhol & 1021 & 197 & 55,91\% & 62,44\% & 58,99\% \\
\textbf{ERMS} & Espanhol & 1057 & 225 & 61,65\% & 56,44\% & 58,93\% \\
\textbf{RRBI} & Hindi & 1046 & 117 & 42,42\% & 47,86\% & 44,98\% \\
\textbf{NJS} & Coreano & 968 & 110 & 42,31\% & 40,00\% & 41,12\% \\
\textbf{ABA} & Árabe & 959 & 82 & 32,31\% & 25,61\% & \textbf{28,57\%} \\
\bottomrule
\end{tabular}
\end{table}

A análise revela uma discrepância notável na capacidade do modelo de identificar desvios entre diferentes grupos de sotaques. O detector obteve excelente performance para falantes com sotaque **vietnamita** (\textbf{TLV}: F1 de 77,45\% e **THV**: F1 de 72,51\%) e **chinês** (\textbf{HQTV}: F1 de 74,29\%). Coincidentemente, estes falantes registram o maior volume absoluto de erros reais no corpus. 

Em contrapartida, falantes de sotaque **árabe** (\textbf{ABA}) apresentaram o menor F1-score (**28,57\%**), com baixas taxas de precisão (32,31\%) e revocação (25,61\%). Esse resultado sugere que os desvios cometidos por falantes árabes são caracterizados por substituições de vogais sutis e variações róticas que o modelo acústico treinado de forma geral tende a classificar como aceitáveis (não disparando erros), ou que a quantidade reduzida de desvios reais nestes falantes distorce o peso relativo dos falsos positivos.

\section{Conclusão}
Este trabalho apresentou o desenvolvimento e a validação sistemática de um pipeline baseado em redes de aprendizado profundo do tipo \textit{Transformers} para a tarefa de detecção de erros de pronúncia em nível de fonema. Ao integrar o Wav2Vec 2.0 (como modelo acústico treinado em CPU via backend MPS) com o OpenAI Whisper (como reconhecedor ASR robusto) e alinhamento baseado na distância de Levenshtein, foi possível construir um sistema completo de diagnóstico de desvios linguísticos.

O ajuste fino local do Wav2Vec 2.0 por 25 épocas provou ser altamente eficaz, atingindo um PER de 17,63\% e viabilizando a inferência acústica com acurácia acima de 82\%. A avaliação em massa do detector completo no conjunto de testes de 720 áudios do corpus L2-ARCTIC produziu métricas globais sólidas (F1-score de 55,79\% e acurácia de 87,02\%), com destaque para a alta sensibilidade na identificação de deleções fonéticas (71,81\% de revocação). As análises estratificadas expuseram a variação de performance do detector em função dos sotaques maternos L1 dos estudantes.

Como trabalhos futuros, propõe-se a substituição do reconhecedor acústico base por arquiteturas multilingues maiores (\textit{wav2vec2-large-xlsr-53}) para melhorar a generalização sobre sotaques mais complexos e a incorporação de informações probabilísticas de posterior de fonemas (GOP) combinadas com modelos linguísticos contextuais para suavizar os alinhamentos espúrios.

\section*{Declaração de Uso de IA Generativa}
Em conformidade com as diretrizes da UFRGS e do regulamento da disciplina, declara-se que o assistente inteligente **Antigravity CLI** (baseado em modelos da família Gemini 3.5) foi utilizado no desenvolvimento deste projeto. O assistente auxiliou nas etapas de: (i) especificação da arquitetura e do roadmap inicial; (ii) geração dos scripts de inferência e de alinhamento estruturado de fonemas; (iii) desenvolvimento do script de avaliação quantitativa automática; e (iv) redação e formatação preliminar deste texto em LaTeX. Todo o código gerado e a análise estatística foram revisados, executados e validados na máquina local pelos autores, que assumem integral responsabilidade pelo conteúdo deste artigo científico.

\bibliographystyle{sbc}
\begin{thebibliography}{99}

\bibitem[Baevski et al. 2020]{baevski2020wav2vec}
Baevski, A., Zhou, H., Mohamed, A., and Auli, M. (2020).
\newblock wav2vec 2.0: A framework for self-supervised learning of speech representations.
\newblock {\em Advances in Neural Information Processing Systems}, 33:12449--12460.

\bibitem[Radford et al. 2023]{radford2023robust}
Radford, A., Kim, J.~W., Xu, T., Brockman, G., McLeavey, C., and Sutskever, I. (2023).
\newblock Robust speech recognition via large-scale weak supervision.
\newblock In {\em International Conference on Machine Learning}, pages 28490--28508. PMLR.

\bibitem[Witt and Young 2000]{witt2000phone}
Witt, S.~M. and Young, S.~J. (2000).
\newblock Phone-level pronunciation scoring and assessment for an interactive English language learning system.
\newblock {\em Speech Communication}, 30(2-3):95--108.

\bibitem[Zhao et al. 2018]{zhao2018l2}
Zhao, G., Sonsaat, S., Silpachai, A., Lim, I., Oney, M., and Chukharev-Hudilainen, E. (2018).
\newblock L2-ARCTIC: A non-native English speech corpus for research in pronunciation assessment.
\newblock In {\em Proc. Interspeech 2018}, pages 1150--1154.

\bibitem[Adsouza 2026a]{techchallenges}
Dsouza, A. (2026a).
\newblock Relatório de Gargalos Técnicos e Lições Aprendidas - Projeto TF-PLN.
\newblock {\em Documentação Técnica interna}. UFRGS.

\bibitem[Adsouza 2026b]{evaluate_detector}
Dsouza, A. (2026b).
\newblock Script de avaliação quantitativa em massa (\texttt{evaluate\_detector.py}).
\newblock {\em Repositório de código-fonte}. UFRGS.

\end{thebibliography}

\end{document}
